A positive integer $m$ is \emph{practical} if every integer $x$ with $1\le x<m$ can be written as a sum of \emph{distinct} positive divisors of $m$.

For a practical $m$, define $h(m)$ to be the smallest integer $h\ge 0$ such that every integer $x$ with $1\le x<m$ can be written as a sum of at most $h$ distinct divisors of $m$.

Questions:

(Q1) Are there infinitely many practical $m$ with
\[
h(m) < (\log\log m)^{O(1)}?
\]
Interpreting the $O(1)$: does there exist a constant $C>0$ and infinitely many practical $m$ such that $h(m)\le (\log\log m)^C$?

(Q2) For factorials, is it true that
\[
h(n!) < n^{o(1)}\ ?
\]
More strongly, is $h(n!)<(\log n)^{O(1)}$?

(Edge case: $m=1$ is practical vacuously; then the set $\{x:1\le x<m\}$ is empty, so the minimal $h(1)$ is $0$.)
